% Source: https://artofproblemsolving.com/community/c3h3635924_amc_10_prep
% Stable identity: aops-topic-3635924-q03
For positive base, a power is $1$ exactly when the base is $1$ or the exponent is $0$. Base $1$ gives $x^2-4x+1=0$, with roots $2\pm\sqrt 3$ and sum $4$. Exponent $0$ gives $x^2-5x+2=0$, with roots $\dfrac{5\pm\sqrt{17}}{2}$ and sum $5$; the base at these roots is $x>0$. A negative base can yield $1$ only if the base is $-1$ and the real-defined rational exponent has even numerator. Base $-1$ occurs at $x=1,3$; the exponents are $-2,-4$, so both qualify. Base $0$ never qualifies. The six distinct roots therefore sum to $4+5+1+3=\boxed{13}$.
